% Generated by Listas/make_statements.py from pratica-lecture-05.md, item 8. Do not edit.
\textbf{Source.} Tadelis (2013), \emph{Game Theory: An Introduction}, Exercise 9.3 ``Campaigning Revisited'', pp.~187--188 {[}PDF 206--207{]}. Answer in the source: none (the book prints no solutions; the preface says about half are on Princeton UP's website, not checked). Reproduced verbatim.



\begin{quote}

\textbf{Campaigning Revisited:} Two political candidates are scheduled to campaign in two states, in one in period \(t = 1\) and in the other in \(t = 2\). In each state each candidate can choose either a positive campaign that promotes his own agenda (\(P\) for player 1, \(p\) for player 2) or a negative one that attacks his opponent (\(N\) for player 1, \(n\) for player 2). Residents of the first-period state do not mind negative campaigns, which are generally effective, and payoffs in this state are given by the following matrix:



{\def\LTcaptype{none} % do not increment counter

\begin{longtable}[]{@{}ccc@{}}

\toprule\noalign{}

Player 1 ~Player 2 & \(p\) & \(n\) \\

\midrule\noalign{}

\endhead

\bottomrule\noalign{}

\endlastfoot

\(P\) & 2, 2 & 0, 5 \\

\(N\) & 5, 0 & 3, 3 \\

\end{longtable}

}



In the second-period state, residents dislike negative campaigns despite their effectiveness, and the payoffs are given by the following matrix:



{\def\LTcaptype{none} % do not increment counter

\begin{longtable}[]{@{}ccc@{}}

\toprule\noalign{}

Player 1 ~Player 2 & \(p\) & \(n\) \\

\midrule\noalign{}

\endhead

\bottomrule\noalign{}

\endlastfoot

\(P\) & 4, 4 & 1, 0 \\

\(N\) & 0, 1 & 2, 2 \\

\end{longtable}

}



\begin{enumerate}

\def\labelenumi{\alph{enumi}.}

\item

  What are the Nash equilibria of each stage-game? Find all the pure-strategy subgame-perfect equilibria with extreme discounting (\(\delta = 0\)). Be precise in defining history-contingent strategies for both players.

\item

  Now let \(\delta = 1\). Find a subgame-perfect equilibrium for the two-stage game in which the players choose \((P, p)\) in the first stage-game.

\item

  What is the lowest value of \(\delta\) for which the subgame-perfect equilibrium you found in (b) survives?

\item

  Can you find a subgame-perfect equilibrium for this game in which the players play something other than \((P, p)\) or \((N, n)\) in the first stage?

\end{enumerate}

\end{quote}



\textbf{Note (not part of the exercise).} As printed, parts (b)--(d) have no solution with \(\delta \le 1\): the best-reply gain from \(N\) in stage 1 is 3, while stage 2 can swing a player's payoff by at most \(4 - 2 = 2\). The solution proves this and gives the threshold \(\delta = 3/2\).

